%%%%%%%%%%%%%%%%%%%%% chapter.tex %%%%%%%%%%%%%%%%%%%%%%%%%%%%%%%%%
%
% sample chapter
%
% Use this file as a template for your own input.
%
%%%%%%%%%%%%%%%%%%%%%%%% Springer Nature %%%%%%%%%%%%%%%%%%%%%%%%%%
\motto{Henceforth space by itself, and time by itself, are doomed to fade 
away into mere shadows, and only a kind of union of the two will preserve 
an independent reality. ---Hermann Minkowski}
\chapter{Spacetime Curvature}
\label{curvature}

\abstract*{
  Since the introduction of the general theory of relativity geometry 
  has become an implicit property of spacetime in humanity's understanding
  of the universe. It stands as the most beautiful and most mystical theory 
  because of its proven predictions. The more that is observed in the universe, 
  the better the theory becomes, which is truly unique. To grasp, what it
  means for a scientific theory to be beautiful, one has to dive deep into 
  the concepts and notions entailed in it. Before this gets a practical and 
  experimental use, the introduction of spacetime curvature is imminent, so 
  it will be covered in this chapter. There will be somewhat rigorous 
  explanations of the purely mathematical concepts but at a level for 
  a physicist's knowledge.
}

\abstract{
  Since the introduction of the general theory of relativity geometry 
  has become an implicit property of spacetime in humanity's understanding
  of the universe. It stands as the most beautiful and most mystical theory 
  because of its proved predictions. The more is observed in the universe, 
  the better the theory becomes, that is truly unique. To grasp, what it
  means for a scientific theory to be beautiful, one has to dive deep into 
  the concepts and notions entailed in it. Before this gets a practical and 
  experimental use, the introduction of spacetime curvature is imminent, so 
  it will be covered in this chapter. There will be somewhat rigorous 
  explanations of the purely mathematical concepts but at a level for 
  a physicist's knowledge.
}

\section{Manifolds}
\label{sec:manifolds}

The whole general theory of relativity is built on the concept of 
manifolds, in particular \emph{Pseudo Riemannian Manifolds}.
One does need to work with these notions because the four dimensional 
euclidean spacetime would not be feasible for the complex geometry 
of the actual spacetime.

\subsection{Definition}
\label{subsec:def-manifold}

Every mathematician is used to euclidean spaces of type $\mathbb{R}^n$ 
with dimensions $n$ where space is the same everywhere and differentiation, 
integration and other trivial operations are properly defined and used 
widely. Clearly, this notion does not apply to the spacetime surrounding 
us. For that a new concept is needed and that concept is a manifold.
A manifold, in general, is any geometric object that does, globally, not 
look like $\mathbb{R}^n$ but does so locally. For instance, any type of 
elliptical object, including the two sphere $S^2$, is a 
manifold. Another condition that manifolds must fulfill is the constant 
dimension $n$ on the whole manifold. So, a line going into a plane is 
not a manifold because it would imply a transition from $\mathbb{R}$ to 
$\mathbb{R}^2$ where the line meets the plane~\cite{Book:general_relativity_03_carroll}.

If one is looking for a precise mathematical definition, a manifold 
of the dimension $n$ is a topological space where every point $P$ has a neighborhood that is homeomorphic to an open subset of $\mathbb{R}^n$.

\paragraph{Coordinate Charts}

A coordinate chart $\phi$ is a projection from a manifold $\mathcal{M}$ of 
a dimension $n$ to a continuous function in $\mathbb{R}^n$ that can 
be expressed as 
\begin{equation}
  \phi: \mathcal{M} \rightarrow \mathbb{R}^n. 
  \label{eq:chart}
\end{equation}



\subsection{Tangent Space}
\label{subsec:tangent-space}



\subsection{Cotangent Space}
\label{subsec:cotangent-space}

\subsection{Coordinate Basis}
\label{subsec:coordinate-basis}

\section{The Metric}
\label{sec:metric}

The metric $g_{\mu \nu}$ is the fundamental abstract object in the general 
theory of relativity. 

\subsection{Geometric Definition}
\label{subsec:geo-def-metric}

\subsection{Mathematical Definition}
\label{subsec:math-def-metric}

\section{Christoffel Symbol}
\label{sec:christoffel}

\section{Ricci Tensor}
\label{sec:ricci-tensor}

\subsection{Commutator}
\label{subsec:commutator}

\subsection{Riemann Tensor}
\label{subsec:riemann-tensor}

\subsection{Ricci Scalar}
\label{subsec:ricci-scalar}

%%%%%%%%%%%%%%%%%%%%%%%%%%%%%%%%%%%%%%%%%%%%%%%%%%%%%%%%%%%%%%%%%%%%%%%%%%%
% References
%%%%%%%%%%%%%%%%%%%%%%%%%%%%%%%%%%%%%%%%%%%%%%%%%%%%%%%%%%%%%%%%%%%%%%%%%%%

\bibliographystyle{spmpsci}
\bibliography{ref-curvature}
